\section{补充实验：把“三轴”搬到没有白盒的任务上}\label{sec:appB}

\noindent\textbf{这一节是主报告的补充，不是替代。}
主报告的结论有一个前提：算法已知。节点分类没有精确算法——没有人能写出“给定图，输出每个节点的类”的正确程序。
这正好是黑盒的主场。于是问题变成：\emph{当白盒缺席时，我们还能不能验证解释是对的？}

\subsection{设定}

\textbf{数据}：Cora（2708 节点、1433 维词袋、7 类、去重后 5278 条边）。Planetoid 公开划分：训练 140（节点 0--139）、验证 500、测试 1000。同配率 0.811。

\textbf{模型}：手写 2 层 GCN（hidden 16、dropout 0.5、$\hat A=D^{-1/2}(A+I)D^{-1/2}$），
以及结构相同的 MLP 作无结构对照——参数量都是 23{,}040，唯一差别是每层是否左乘 $\hat A$。

\textbf{解释器}：边掩码（GNNExplainer 简化版），冻结模型，学掩码 $m$ 最小化
$-\log p(c\mid A\odot m)+\lambda\,\mathrm{mean}(m)$（$\lambda=0.05$，200 步，温度固定 1）。

\textbf{真值}：Cora 上没有“正确的解释”，所以另造合成 SBM（标签由社区决定），
用\textbf{消融}定义每条边的真实重要性 $\Delta_e$：删掉它，目标类概率下降多少。

\subsection{结果}

\begin{table}[H]
\centering
\caption{B 线主要数字}
\label{tab:appB}
\small
\setlength{\tabcolsep}{3pt}
\begin{tabular}{lll}
\toprule
维度 & 结果 & 说明 \\
\midrule
节点分类 & GCN 0.8162，MLP 0.5539 & 差 26 个点，只能归因于图结构 \\
随机加边 10\% & $-0.017$ & 对随机扰动相当耐受 \\
FGA 10\% & $-0.119$ & 同预算约为随机的 7 倍 \\
\midrule
SBM：解释边 $\Delta$ & 0.0673 & 随机边 0.0044（Wilcoxon $p=5\times10^{-11}$） \\
SBM：oracle 效率 & 中位 0.939 & 达到“按真值排序”的 94\%（41/44 节点可算） \\
SBM：无法判定 & 1/45 & 掩码并列，排不出序 \\
\midrule
Cora：掩码可判定率 & 88\%（105/120） & 候选边中位 48 条仍能给出排序 \\
Cora：充分性 & 1.31 vs 随机 1.21 & 保留解释边后目标类概率 \\
Cora：必要性 & 删解释边 $-0.123$ vs 删随机边 $-0.019$ & 解释边确实更“必要” \\
Cora：跨模型 Jaccard & 0.714 & 换训练种子后解释基本一致 \\
\midrule
攻击后解释 & Jaccard 1.000 & \textbf{解释对 FGA 完全不敏感}（见下） \\
负控制 & 3/3 & bottom-$k$ 桩实现把 $\Delta$ 压到 0.08 倍 \\
\bottomrule
\end{tabular}
\end{table}

\subsection{两条最有意思的发现}

\textbf{第一，解释能被验证，但必须借合成环境。}在 SBM 上我们造出了真值（消融 $\Delta$），
解释器 top-$k$ 的平均重要性达到 oracle 的 94\%，而随机边只有 oracle 的 6\%。
换句话说，\emph{没有白盒时，“解释对不对”仍然可测——但需要一个你知道答案的 Synthetic 环境}。
在真实数据上只能测保真度与稳定性，而这两者都只是必要条件。

\textbf{第二，攻击不改变解释。}用 FGA 加 5\% 的边使准确率下降 7.5 个点，
但解释边集合与攻击前的中位 Jaccard 是 \textbf{1.000}——几乎逐条相同。
这说明解释器给出的边反映的是\emph{局部结构本身}（哪些邻居对目标节点的表示重要），
而不是模型在攻击后“改变的主意”。对可解释性而言这是双刃：
好处是解释稳定、不受攻击噪声影响；坏处是它也不会告诉你“攻击改动了哪里”。

\subsection{三条限制（都必须写）}

\begin{enumerate}[leftmargin=1.6em]
\item \textbf{解释器实现上踩过一个致命 bug，已修。}第一版归一化算子把度数重复累加、
自环写成 $d$ 而非 $d^2$，且只传单方向边表，导致掩码塌成常数、可判定率 0/120——那段时间的所有解释结论都是错的。
修复后与训练算子逐元素一致（全图 max 差 0.0，删边后 $3\times10^{-8}$）。
\item \textbf{SBM 上仍有 1/45 个节点排不出序}，且候选池大小与度数分布未落盘，跳过是否有偏无法审计。
\item \textbf{只测了 $\varepsilon=5\%$ 一个攻击预算}，未做跨预算的解释稳定性曲线；也未测 $k$ 的敏感性（$k=6$ 是主报告沿用的预算）。
\end{enumerate}

\subsection{与主报告的关系}

主报告证明：算法已知时，白盒在准确率、鲁棒性、可解释性上同时占优，且保真度会被常数预测器骗过。
本节把同一套怀疑搬到白盒缺席的任务上，结论是：

\begin{itemize}[leftmargin=1.6em]
\item 解释的\emph{验证}依然可行，但代价是要先造一个带真值的合成环境；
\item 在真实数据上，保真度与稳定性都能被满足，却都只是必要条件——
攻击完全不动解释，就说明解释未必跟踪模型当前的证据；
\item 因此“可解释性”在黑盒场景下更像一种\emph{审计工具}：它能告诉你模型依赖哪些局部结构，
但不能替代“误差有没有上界”这种白盒式的保证。
\end{itemize}

这与主报告的结论并不矛盾：白盒缺席之后，连“验证”本身都变得更贵、更间接。
